\chapter{Methodology, Data \& Observations}
\label{sec:methodology-chapter}

\section{Methodology}
\label{sec:methodology}

This section states the design principles adopted \emph{before} any model
was trained or any result existed, the resulting architecture, and the
validation protocol used to judge every result reported in
Section~\ref{sec:results}. Table~\ref{tab:decisionlog} at the end of this
section is a decision log: for each major design choice, the alternatives
considered, the choice made, and the evidence that settled it. Most write-ups
of a pipeline like this present only the final configuration; the
alternatives that were rejected, and why, are at least as informative.

\subsection{Design principles}

\paragraph{Wavelength conditioning, not band identity.}
ZTF alerts carry three photometric filters; LSST alerts carry six
(Table~\ref{tab:lameff}). If a model's input encoded ``filter index 2,''
ZTF's second filter and LSST's second filter would be unrelated integers
occupying unrelated positions in the model's input space, and nothing
learned from simulated LSST photometry (Tier A) would transfer to real ZTF
photometry (Tier B) at all. Instead, every observation in this pipeline
carries its filter's actual \emph{effective wavelength}, a continuous
physical quantity, linearly mapped into $[-1, 1]$
(Section~\ref{sec:analysis}). ZTF's $r$ band (633.9\,nm) and LSST's $r$ band
(622.2\,nm) land near each other in the model's input space automatically,
because they are physically similar wavelengths, not because any code
special-cases ``these are both called $r$.'' This single decision is what
makes ``pretrain on simulated LSST photometry, fine-tune on real ZTF
photometry'' a coherent strategy rather than two unrelated models sharing a
name.

\begin{table}[htb]
\centering
\caption{Filter effective wavelengths used for wavelength conditioning
(Section~\ref{sec:methodology}). LSST's six bands and ZTF's three are
mapped into the same continuous space; no band index is ever used as a
model input.}
\label{tab:lameff}
\begin{tabular}{ccccc}
\toprule
\multicolumn{2}{c}{LSST} & & \multicolumn{2}{c}{ZTF} \\
Band & $\lambda_\mathrm{eff}$ (nm) & & Band & $\lambda_\mathrm{eff}$ (nm) \\
\midrule
$u$ & 368.6 & & $g$ & 472.3 \\
$g$ & 480.3 & & $r$ & 633.9 \\
$r$ & 622.2 & & $i$ & 788.6 \\
$i$ & 754.5 & & & \\
$z$ & 869.1 & & & \\
$y$ & 971.0 & & & \\
\bottomrule
\end{tabular}
\end{table}

\paragraph{Causal models only.}
Every recurrent layer in this pipeline runs in a single direction only
(\texttt{bidirectional=False} throughout), with no attention mechanism
that could attend to a future timestep. A bidirectional model evaluated
offline, where
every object's complete light curve already exists in the training set,
would silently learn to use information from an object's future relative to
whatever point it is nominally classifying at -- on day 3 of a real alert
stream, day 60 has not happened yet. Such a model scores well offline and is
meaningless in production, and the failure would not be visible in any
ordinary train/test split, since both sides of that split share the same
non-causal leakage. This is treated as an absolute architectural constraint,
not a tunable hyperparameter, and is enforced by a dedicated unit test that
truncates a real input sequence at a timestep $t$ and asserts the model's
prediction at $t$ is bit-identical whether or not timesteps after $t$ are
present in the input tensor at all (Table~\ref{tab:gates}, Gate B).

\paragraph{A frozen evaluation harness.}
The canonical data schema and the evaluation metrics code were written and
frozen before any model existed, and were never modified afterward except
by one documented, additive amendment (below). The reasoning is that a
metric implementation that can still be edited is a metric that will,
eventually, be adjusted -- consciously or not -- to make a disappointing
result look better, and the resulting number would be unfalsifiable: there
would be no way to tell ``the model improved'' from ``the yardstick moved''
after the fact. The one amendment made to the frozen harness after Day~1
was the addition of ROC-AUC (Ia-vs-rest) as a new metric, once the project's
own literature anchors (Section~\ref{sec:literature}) turned out to name a
metric that did not yet exist anywhere in the harness. The amendment added a
new function and the minimal wiring to carry its output through the existing
report structures, touched no existing function's computed output, and was
verified independently (a new value on an existing checkpoint, not a
retroactive recomputation of an existing reported number). An amendment of
this kind -- additive, independently verified, motivated by a real gap
rather than a disappointing result -- is judged preferable to a blanket
no-changes rule that would otherwise have blocked a metric the project's
own anchors required.

\paragraph{Flux, never magnitude.}
Difference-image photometry -- subtracting a reference template from a new
science exposure -- routinely produces negative flux values on noisy,
faint, non-detections. Those points are concentrated in exactly the early,
pre-peak regime this project's early-classification goal cares about most.
A magnitude is a logarithm of flux and cannot represent a negative value at
all; any reader that converts to magnitude before this pipeline sees the
data has already silently discarded some of the most informative early-time
points. Every reader in this pipeline therefore emits flux, on one shared
zeropoint (Section~\ref{sec:data}), and only ever converts to magnitude for
plotting.

\paragraph{A canonical alert schema as the single conversion target.}
Every data source -- simulated ELAsTiCC2, real ZTF/BTSbot, live Fink -- is
converted into one shared \texttt{CanonicalAlert} record before any model
or evaluation code ever touches it (Table~\ref{tab:schema}). This is what
makes ``pretrain on one source, fine-tune on another, serve a third''
possible without a source-specific branch anywhere downstream of the reader
layer.

\begin{table}[htb]
\centering
\caption{The canonical alert schema (\texttt{src/schema.py}, frozen after
Day 1--2 of the project). Every reader converts its source into this record;
no downstream code depends on which survey an alert came from.}
\label{tab:schema}
\begin{tabular}{L{2.6cm}L{10.9cm}}
\toprule
Field & Meaning \\
\midrule
\texttt{object\_id} & Survey-native identifier \\
\texttt{survey} & ZTF / LSST / ELAsTiCC \\
\texttt{observations} & List of (mjd, $\lambda_\mathrm{eff}$, flux, flux\_err, is\_detection), sorted ascending by mjd \\
\texttt{static} & Host photo-$z$, offset, ellipticity, host magnitude, MW $E(B\!-\!V)$, galactic latitude, real-bogus score \\
\texttt{cutouts} & $(3,63,63)$ science/reference/difference postage stamps, or \texttt{None} \\
\texttt{class\_id} & ELAsTiCC2 taxonomy id (Table~\ref{tab:taxonomy}), or \texttt{None} if unlabelled \\
\bottomrule
\end{tabular}
\end{table}

\subsection{Figure 1: pipeline architecture}

Figure~\ref{fig:pipeline} is the anchor figure for this report: the three
data tiers on the left, the canonical schema as the single conversion
funnel every tier passes through, the shared light-curve encoder feeding
both the classification head and the anomaly scorer, and the serving path
from a live alert stream through to a results interface.

\begin{figure}[htb]
\centering
\resizebox{\linewidth}{!}{% Figure 1: pipeline architecture. Included via % Figure 1: pipeline architecture. Included via % Figure 1: pipeline architecture. Included via \input{figures/fig1_pipeline}
% from paper/sections/methodology.tex.
\begin{tikzpicture}[
  font=\small,
  box/.style={draw, rounded corners, align=center, minimum height=0.9cm, inner sep=5pt},
  tier/.style={box, minimum width=3.4cm, fill=gray!8},
  funnel/.style={draw, rounded corners, minimum width=2.9cm, minimum height=2.6cm, align=center, fill=gray!15},
  enc/.style={box, minimum width=4.8cm, fill=blue!6},
  head/.style={box, minimum width=3.6cm},
  serve/.style={box, minimum width=2.0cm, fill=gray!5, font=\scriptsize},
  arr/.style={-{Stealth[length=2.2mm]}, thick},
  >=Stealth,
]

% --- data tiers (left column) ---
\node[tier] (tiera) at (0, 5.4) {Tier A\\ELAsTiCC2 (simulated)\\843{,}617 objects};
\node[tier] (tierb) at (0, 3.6) {Tier B\\BTSbot / ZTF BTS (real)\\19{,}511 objects};
\node[tier] (tierc) at (0, 1.8) {Tier C\\Fink live stream (real)\\unlabelled};

% --- canonical schema funnel ---
\node[funnel] (schema) at (4.3, 3.6) {\textbf{CanonicalAlert}\\(frozen schema)};

\draw[arr] (tiera.east) -- (schema.north west);
\draw[arr] (tierb.east) -- (schema.west);
\draw[arr] (tierc.east) -- (schema.south west);

% --- encoder ---
\node[enc] (encoder) at (9.7, 3.6) {Causal GRU light-curve encoder\\(+ optional static host branch)};
\draw[arr] (schema) -- (encoder);

% --- heads ---
\node[head, draw=classIa, thick] (clshead) at (9.7, 6.0) {3-class head\\\textcolor{classIa}{Ia} / \textcolor{classIbc}{Ib/c} / \textcolor{classII}{II}};
\node[head] (anomhead) at (9.7, 1.2) {Anomaly scorer\\recon.\ error + Mahalanobis + iso.\ forest};
\draw[arr] (encoder.north) -- (clshead.south);
\draw[arr] (encoder.south) -- (anomhead.north);

% --- serving path (bottom row) ---
\node[serve] (kafka) at (0.0, -1.4) {Kafka};
\node[serve] (redis) at (2.15, -1.4) {Redis\\rolling state};
\node[serve] (onnx) at (4.3, -1.4) {ONNX\\inference};
\node[serve] (pg) at (6.45, -1.4) {Postgres};
\node[serve] (api) at (8.6, -1.4) {API /\\dashboard};

\draw[arr] (kafka) -- (redis);
\draw[arr] (redis) -- (onnx);
\draw[arr] (onnx) -- (pg);
\draw[arr] (pg) -- (api);

\draw[arr] (tierc.south) .. controls +(0,-1.6) and +(-2.2,1.0) .. (kafka.north)
  node[midway, below left, font=\scriptsize] {replay or live};
\draw[arr] (clshead.south east) .. controls +(1.2,-2.0) and +(1.0,3.4) .. (onnx.north east);
\draw[arr] (anomhead.south east) .. controls +(0.8,-0.6) and +(1.4,0.9) .. (onnx.north east);

\end{tikzpicture}

% from paper/sections/methodology.tex.
\begin{tikzpicture}[
  font=\small,
  box/.style={draw, rounded corners, align=center, minimum height=0.9cm, inner sep=5pt},
  tier/.style={box, minimum width=3.4cm, fill=gray!8},
  funnel/.style={draw, rounded corners, minimum width=2.9cm, minimum height=2.6cm, align=center, fill=gray!15},
  enc/.style={box, minimum width=4.8cm, fill=blue!6},
  head/.style={box, minimum width=3.6cm},
  serve/.style={box, minimum width=2.0cm, fill=gray!5, font=\scriptsize},
  arr/.style={-{Stealth[length=2.2mm]}, thick},
  >=Stealth,
]

% --- data tiers (left column) ---
\node[tier] (tiera) at (0, 5.4) {Tier A\\ELAsTiCC2 (simulated)\\843{,}617 objects};
\node[tier] (tierb) at (0, 3.6) {Tier B\\BTSbot / ZTF BTS (real)\\19{,}511 objects};
\node[tier] (tierc) at (0, 1.8) {Tier C\\Fink live stream (real)\\unlabelled};

% --- canonical schema funnel ---
\node[funnel] (schema) at (4.3, 3.6) {\textbf{CanonicalAlert}\\(frozen schema)};

\draw[arr] (tiera.east) -- (schema.north west);
\draw[arr] (tierb.east) -- (schema.west);
\draw[arr] (tierc.east) -- (schema.south west);

% --- encoder ---
\node[enc] (encoder) at (9.7, 3.6) {Causal GRU light-curve encoder\\(+ optional static host branch)};
\draw[arr] (schema) -- (encoder);

% --- heads ---
\node[head, draw=classIa, thick] (clshead) at (9.7, 6.0) {3-class head\\\textcolor{classIa}{Ia} / \textcolor{classIbc}{Ib/c} / \textcolor{classII}{II}};
\node[head] (anomhead) at (9.7, 1.2) {Anomaly scorer\\recon.\ error + Mahalanobis + iso.\ forest};
\draw[arr] (encoder.north) -- (clshead.south);
\draw[arr] (encoder.south) -- (anomhead.north);

% --- serving path (bottom row) ---
\node[serve] (kafka) at (0.0, -1.4) {Kafka};
\node[serve] (redis) at (2.15, -1.4) {Redis\\rolling state};
\node[serve] (onnx) at (4.3, -1.4) {ONNX\\inference};
\node[serve] (pg) at (6.45, -1.4) {Postgres};
\node[serve] (api) at (8.6, -1.4) {API /\\dashboard};

\draw[arr] (kafka) -- (redis);
\draw[arr] (redis) -- (onnx);
\draw[arr] (onnx) -- (pg);
\draw[arr] (pg) -- (api);

\draw[arr] (tierc.south) .. controls +(0,-1.6) and +(-2.2,1.0) .. (kafka.north)
  node[midway, below left, font=\scriptsize] {replay or live};
\draw[arr] (clshead.south east) .. controls +(1.2,-2.0) and +(1.0,3.4) .. (onnx.north east);
\draw[arr] (anomhead.south east) .. controls +(0.8,-0.6) and +(1.4,0.9) .. (onnx.north east);

\end{tikzpicture}

% from paper/sections/methodology.tex.
\begin{tikzpicture}[
  font=\small,
  box/.style={draw, rounded corners, align=center, minimum height=0.9cm, inner sep=5pt},
  tier/.style={box, minimum width=3.4cm, fill=gray!8},
  funnel/.style={draw, rounded corners, minimum width=2.9cm, minimum height=2.6cm, align=center, fill=gray!15},
  enc/.style={box, minimum width=4.8cm, fill=blue!6},
  head/.style={box, minimum width=3.6cm},
  serve/.style={box, minimum width=2.0cm, fill=gray!5, font=\scriptsize},
  arr/.style={-{Stealth[length=2.2mm]}, thick},
  >=Stealth,
]

% --- data tiers (left column) ---
\node[tier] (tiera) at (0, 5.4) {Tier A\\ELAsTiCC2 (simulated)\\843{,}617 objects};
\node[tier] (tierb) at (0, 3.6) {Tier B\\BTSbot / ZTF BTS (real)\\19{,}511 objects};
\node[tier] (tierc) at (0, 1.8) {Tier C\\Fink live stream (real)\\unlabelled};

% --- canonical schema funnel ---
\node[funnel] (schema) at (4.3, 3.6) {\textbf{CanonicalAlert}\\(frozen schema)};

\draw[arr] (tiera.east) -- (schema.north west);
\draw[arr] (tierb.east) -- (schema.west);
\draw[arr] (tierc.east) -- (schema.south west);

% --- encoder ---
\node[enc] (encoder) at (9.7, 3.6) {Causal GRU light-curve encoder\\(+ optional static host branch)};
\draw[arr] (schema) -- (encoder);

% --- heads ---
\node[head, draw=classIa, thick] (clshead) at (9.7, 6.0) {3-class head\\\textcolor{classIa}{Ia} / \textcolor{classIbc}{Ib/c} / \textcolor{classII}{II}};
\node[head] (anomhead) at (9.7, 1.2) {Anomaly scorer\\recon.\ error + Mahalanobis + iso.\ forest};
\draw[arr] (encoder.north) -- (clshead.south);
\draw[arr] (encoder.south) -- (anomhead.north);

% --- serving path (bottom row) ---
\node[serve] (kafka) at (0.0, -1.4) {Kafka};
\node[serve] (redis) at (2.15, -1.4) {Redis\\rolling state};
\node[serve] (onnx) at (4.3, -1.4) {ONNX\\inference};
\node[serve] (pg) at (6.45, -1.4) {Postgres};
\node[serve] (api) at (8.6, -1.4) {API /\\dashboard};

\draw[arr] (kafka) -- (redis);
\draw[arr] (redis) -- (onnx);
\draw[arr] (onnx) -- (pg);
\draw[arr] (pg) -- (api);

\draw[arr] (tierc.south) .. controls +(0,-1.6) and +(-2.2,1.0) .. (kafka.north)
  node[midway, below left, font=\scriptsize] {replay or live};
\draw[arr] (clshead.south east) .. controls +(1.2,-2.0) and +(1.0,3.4) .. (onnx.north east);
\draw[arr] (anomhead.south east) .. controls +(0.8,-0.6) and +(1.4,0.9) .. (onnx.north east);

\end{tikzpicture}
}
\caption{Pipeline architecture. Three data tiers (simulated ELAsTiCC2 for
pretraining, real ZTF/BTSbot for fine-tuning, live Fink for serving) convert
into one canonical alert schema, which every downstream stage consumes
identically regardless of source. A shared causal GRU light-curve encoder
(optionally fused with a static host-context branch) feeds both the flat
3-class classification head and a separate anomaly-scoring branch. The
serving path (Kafka $\to$ Redis rolling state $\to$ ONNX inference $\to$
Postgres $\to$ API) is the same consumer code whether its producer is an
archived-data replay or the live Fink stream.}
\label{fig:pipeline}
\end{figure}

\subsection{Model architecture}

The light-curve encoder (Figure~\ref{fig:pipeline}) is a 2-layer causal GRU,
hidden size 128, with a \texttt{LayerNorm} and dropout ($p=0.2$) applied to
the GRU's output before a linear head. Sequences are packed and masked
(\texttt{pack\_padded\_sequence}) so that a padded timestep is never
processed by the recurrent layer for any sequence in a batch, regardless of
where it sits in the batch. Predictions are emitted at \emph{every}
timestep, not only the last, which is what makes the phase-cutoff curves in
Section~\ref{sec:results} possible without a separate model per cutoff.

Each observation is represented by six per-timestep features:
$\Delta t$ (days since the object's own first detection),
$\mathrm{signed\_log1p}(\mathrm{flux}/\mathrm{flux\_err})$ (a compressed
signal-to-noise ratio), $\mathrm{signed\_log1p}(\mathrm{flux})$ (compressed
absolute flux), $\log1p(\mathrm{flux\_err})$, the linearly-mapped filter
wavelength, and a detection/forced-photometry indicator. Absolute flux is
carried \emph{in addition to} the signal-to-noise ratio, not instead of it,
even though it is in principle recoverable as their product
($\mathrm{flux} = (\mathrm{flux}/\mathrm{flux\_err}) \times
\mathrm{flux\_err}$): a Type Ia supernova's absolute peak brightness clusters
tightly near $M \approx -19.3$, while Ib/c and II are both fainter and more
scattered, so absolute brightness is real discriminating signal for exactly
the Ia/non-Ia boundary this pipeline is most sensitive on -- but that signal
is only accessible as a \emph{multiplicative} interaction between two other,
already nonlinearly-compressed input channels, and GRUs and MLPs are
comparatively poor at learning multiplicative interactions between their own
inputs. Supplying $\mathrm{signed\_log1p}(\mathrm{flux})$ directly removes
that burden at no additional leakage cost (Section~\ref{sec:analysis}).

An optional static-context branch (a small MLP over seven host/extinction
fields) is fused with the GRU's output before the head in the
``LC+static'' configuration (Section~\ref{sec:results}). Missingness in
these fields is handled explicitly: each field carries its own
zero-filled, z-scored value \emph{and} a separate missingness indicator,
rather than a silently imputed value alone. A zero-fill on already
z-scored data lands at the training-set mean, a neutral placeholder in the
untransformed scale as well; the indicator lets the downstream MLP learn to
treat a filled value differently from a measured one, which a silent
imputation would not allow it to do at all.

\subsection{Anomaly branch}

The anomaly score combines three components, computed over the encoder's
own latent representation and the observed light curve:
\begin{itemize}
  \item \textbf{Reconstruction error}: a small decoder predicts each
    timestep's flux-derived features from the encoder's hidden state as of
    the \emph{previous} timestep only (never the current one, which the GRU
    has already consumed and would make reconstruction trivial), trained on
    the three normal classes only. An object whose flux behaviour this
    decoder cannot predict well is anomalous by construction.
  \item \textbf{Mahalanobis distance} of an object's latent encoding from a
    Gaussian fitted to the normal-class population's latent space.
  \item \textbf{Isolation forest} (\texttt{coniferest}'s \texttt{PineForest},
    used here in pure unsupervised mode) over the concatenation of the
    latent encoding and a handful of directly-engineered light-curve summary
    statistics.
\end{itemize}
The three components are z-scored against their own fit-population
statistics and combined by an equal-weighted average
(Section~\ref{sec:results} discusses why entropy, a fourth candidate
component, was tried and dropped).

\subsection{Hierarchical loss (WHXE)}

Weighted hierarchical cross-entropy \citep{villar2023whxe} was implemented
and compared, in one controlled run (same encoder, data, split, seed, and
class weighting as the flat-head baseline), against the flat 3-class head
used everywhere else in this report. On the data available at the time of
that comparison -- every training object was, by construction, one of the
three target SN-like leaves, with no non-SN-like sibling for the hierarchy
to correctly reject -- WHXE scored \emph{worse} than the flat head
(macro-F1 \TierWHXEMacroFFull{} versus \TierLCOnlyMacroFFull{}, a
$\WHXEMacroFDelta{}$ difference), because the loss's only lever, redistributing
confidence toward a shared parent node, had nowhere useful to spend that
redistribution on a dataset with a single populated branch. Node-level
accuracy at the ``SN-like'' parent measured $\WHXENodeAccuracy{}$ on this
same run, but that number is not informative on its own: every object in
this dataset \emph{is} SN-like, so there is no wrong branch for the
hierarchy to have been tested against. This report treats WHXE as
provisional, not rejected: its entire justification is graceful degradation
in the presence of genuine non-supernova contaminants, a property this
particular comparison could not test, since no such contaminant existed in
the data available for it.

A genuine ambiguity in the published loss's own notation was found and
resolved during implementation, and is recorded here because it changes the
loss's qualitative behaviour depending on which reading is taken. The
published loss (their Eq.~2) indexes a class's position in the taxonomy tree
as $c^{(0)}$ at the \emph{leaf}, increasing toward the root, and weights each
level's term by $\lambda(c^{(h)}) = \exp(-\alpha \, h)$ (their Eq.~5). Read
literally with their own leaf-indexed $h$, this weights \emph{fine-grained}
distinctions (near the leaf, where $h$ is small and $\lambda$ near 1) most
heavily and \emph{coarse} distinctions (near the root, where $h$ is large and
$\lambda$ small) least -- the opposite of the paper's own surrounding
prose, which states that a \emph{larger} $\alpha$ should favour
\emph{parent}-level (coarse) classification more strongly. A hierarchical
loss built to weight fine-grained mistakes most and coarse mistakes least
has no purpose over a plain flat softmax, since the entire motivation for a
hierarchical loss is to penalise a coarse mistake (supernova versus
artifact) more than a fine one (Ia versus Ib/c). This implementation
therefore indexes depth \emph{from the root} instead, so that $\lambda$ stays
near~1 for coarse, near-root terms and shrinks for fine-grained, leaf-level
terms as $\alpha$ grows -- consistent with both the paper's own prose and
with this project's taxonomy's stated purpose. This is recorded as a
resolved implementation decision, not an unresolved ambiguity, but is
reported explicitly rather than silently, since a reader implementing the
same loss from the same equation could reasonably make the opposite choice.

\subsection{Serving architecture}

The serving path (Figure~\ref{fig:pipeline}) is Kafka
$\to$ a stateless normalisation step $\to$ Redis (rolling per-object
observation history) $\to$ an ONNX-exported classifier plus the anomaly
scorer $\to$ Postgres (durable prediction storage) $\to$ a FastAPI read
layer. Replay mode (an archived-data producer onto the same Kafka topic) was
built and validated before live mode, so that the live-mode step was a
producer swap plus new reader modules, not a rewrite of the consumer -- the
same consumer code was verified, unmodified, against both an archived-data
replay and real Fink alert traffic (Section~\ref{sec:results}).

\subsection{Deviation from the specified toolchain}

The programme brief specified the Fink client, TensorFlow, and LightGBM.
The Fink client was used as specified, against real ZTF and LSST broker
topics. TensorFlow was not used; PyTorch was used instead. LightGBM was not
used at all. This deviation is stated plainly rather than left for a reader
to notice: an ONNX export of a PyTorch checkpoint gave a clean, framework-
agnostic serving path with no retraining required to change inference
runtimes; the astronomical ecosystem this project depends on most directly
(\texttt{astropy}, \texttt{sncosmo}, \texttt{coniferest}) is more naturally
composed with PyTorch's eager-mode Python interface than with a static
TensorFlow graph; and LightGBM is a gradient-boosted-tree method suited to a
pipeline built on hand-engineered per-object features -- the design ALeRCE's
own balanced hierarchical random forest follows \citep{forster2021alerce}
-- rather than to the raw, per-timestep sequence input the brief's own
recurrent-network framing specifies. Substituting a gradient-boosted-tree
method for a sequence model would have been a different architecture
entirely, not an equivalent implementation choice.

\subsection{Validation protocol}

Four gates (Table~\ref{tab:gates}) were defined \emph{before} the results
they judge, as pre-registered acceptance criteria rather than criteria
chosen to match whatever result was obtained.

\begin{table}[htb]
\centering
\caption{The four validation gates, defined before the results they judge.}
\label{tab:gates}
\small
\begin{tabular}{L{1.5cm}L{4.6cm}L{4.6cm}L{3.6cm}}
\toprule
Gate & What it checks & Why this failure mode matters & Outcome \\
\midrule
A (readers) & Light curves are visually plausible per class; negative
  fluxes are present in the parsed data; real BTS class counts are near
  their published ratio & A reader that runs without crashing can still
  silently produce garbage; only a direct check catches it & Passed \\
B (first model) & Shuffled-label test collapses to chance; truncating a
  sequence never changes an earlier prediction; Ia-vs-rest ROC-AUC sits in
  a pre-registered band & Catches label leakage, non-causal architecture,
  and results that are too good to be true, respectively & Passed (at a
  corrected sample size for the shuffle test) \\
C (fine-tune) & Accuracy measurably drops from simulated to real data & No
  drop at all would indicate a broken, leaking evaluation split, not a
  better model & Passed \\
D (anomaly branch) & Manual inspection of the top-ranked anomalies expects
  them artifact-dominated, per a published real-bogus-filtering anchor &
  Checks whether real-bogus filtering is suppressing genuine rare signal &
  Deferred -- see Section~\ref{sec:results} \\
\bottomrule
\end{tabular}
\end{table}

\begin{table}[htb]
\centering
\caption{Decision log: major design choices, the alternatives considered,
the choice made, and the evidence that settled it. Presenting the rejected
alternatives alongside the final configuration, rather than only the final
configuration, is deliberate.}
\label{tab:decisionlog}
\scriptsize
\begin{tabular}{L{2.1cm}L{2.6cm}L{2.5cm}L{3.5cm}L{3.0cm}}
\toprule
Decision & Options considered & Choice made & Rationale & Evidence \\
\midrule
Band representation & Categorical band index; continuous wavelength &
  Continuous wavelength & ZTF/LSST bands must land in a shared input space
  for transfer to work at all & Design argument, not an ablation (transfer
  would be structurally impossible otherwise) \\
Recurrence direction & Bidirectional; causal (single-direction) & Causal &
  A bidirectional model sees future timesteps during offline training that
  do not exist at serving time & Gate B causality test, passed \\
Photometry scale & Magnitude; flux & Flux, never clipped & Magnitude cannot
  represent the negative fluxes routine in faint, early-time
  difference-image photometry & Negative-flux fraction confirmed present
  in parsed Tier A data (\TierANegFluxMin{}--\TierANegFluxMax{}\%
  per model directory) \\
Evaluation harness & Mutable; frozen with additive-only amendments &
  Frozen, one additive amendment (ROC-AUC) & A mutable metric can be
  adjusted to fit a disappointing result after the fact & One amendment,
  independently verified, no existing function's output changed \\
Left-edge window & Full archived history; 30-day pre-detection lookback &
  30-day lookback & Training data must match what a live alert stream
  actually accumulates by an object's own last detection &
  \WindowMeanDetections{} detections vs.\ \WindowMeanForcedPhot{} forced-
  photometry points/object measured; \WindowFracRemoved{}\% of a sampled
  sequence's observations fall outside the window \\
Right-edge cap & Uncapped; fixed span cap & 200-day cap
  (\texttt{max\_span\_days}) & A long-lived recurring source's tensorized
  span is otherwise unbounded and far outside the training distribution &
  Span ablation: \SpanTwoHundredMacroF{} at 200~d vs.\ \SpanUncappedMacroF{}
  uncapped, $<$1\% of any class affected \\
Hierarchical loss & WHXE; flat head & Flat head (WHXE kept, provisional) &
  WHXE was worse on a single-populated-branch dataset with no non-SN-like
  negatives yet to gracefully degrade against & Controlled comparison:
  \TierWHXEMacroFFull{} (WHXE) vs.\ \TierLCOnlyMacroFFull{} (flat) \\
Anomaly ensemble & With classifier entropy; without & Without (3-component) &
  Entropy inverted under domain shift (Section~\ref{sec:results}) and
  actively hurt the combination on both evaluation sets & Combined score
  improved from \AnomalyTierACombined{} to
  \AnomalyTierACombinedThreeRecomputed{} (Tier A) and \AnomalyTierBCombined{}
  to \AnomalyTierBCombinedThreeRecomputed{} (Tier B) on dropping it \\
Tier B host features & LC+static; LC-only & LC-only (deployed) & Host
  features helped on Tier A but measurably hurt the Tier B fine-tune &
  \LCOnlyMacroF{} (LC-only) vs.\ \WithStaticMacroF{} (LC+static) on
  Tier B test \\
Image encoder & Build; cut & Cut & No remaining schedule day, and Tier B
  cutout data was never loaded across three sessions & Descope ladder
  accounting (Section~\ref{sec:discussion}) \\
\bottomrule
\end{tabular}
\end{table}
